\chapter{Dépendances logiques et certificats de cohérence}

\section{Chaîne causale des certificats : états finis \(\to\) restrictions \(\to\) survie \(\to\) section infinie \(\to\) réalisations}

La chaîne logique de l'intégration ne part pas d'une conclusion archimédienne.
Elle s'organise dans l'ordre suivant :
\[
\begin{tikzpicture}[baseline=(current bounding box.center),>=Latex,
 node distance=7mm and 8mm,
 n/.style={rounded corners,draw=hmtblue,fill=hmtlightblue,
 minimum height=8mm,align=center,font=\small}]
 \node[n] (finite) {estados finitos\\tipados};
 \node[n,right=of finite] (maps) {restricciones\\compatibles};
 \node[n,right=of maps] (surv) {supervivencia\\y poda};
 \node[n,right=of surv] (lim) {sección\\infinita};
 \node[n,right=of lim] (read) {lectores\\y realizaciones};
 \draw[->,thick] (finite)--(maps);
 \draw[->,thick] (maps)--(surv);
 \draw[->,thick] (surv)--(lim);
 \draw[->,thick] (lim)--(read);
\end{tikzpicture}
\]
Un certificat de chiffres vérifie une restriction de la dernière flèche ; il ne
crée pas rétrospectivement les états et ne remplace pas le théorème d'existence
de la limite.

\section{Identités arithmétiques exactes des certificats finis}

Les recensements qui accompagnent le certificat de filtration nonadique satisfont
\[
 396\cdot6=2376,
 \qquad
 43\cdot9=387,
 \qquad
 432+4\cdot144=1008.
\]
La première égalité relie les événements élémentaires de frontière aux trits
émis par canal ; la deuxième compte les paires de phase \((K,K+9)\) ; la
troisième reconstruit la multiplicité totale des cinq régions de
\(\pi\).

Les deux rôles régionaux qui ne peuvent pas être permutés sont
\[
\begin{array}{rcl}
U_\pi^\perp&=&(501,614,498,169,272,272),\\
E_{\rm sig}&=&555555,\qquad \text{multiplicité }432,\\[.35em]
U_\pi^{--}&=&(870,923,810,418,674,521),\\
E_{\rm sig}&=&933717,\qquad \text{multiplicité }144.
\end{array}
\]

\section{Existence et unicité de la branche infinie dans un arbre à branchement fini de préfixes survivants}

\begin{lemma}[Arbre à branchement fini]
Soit \(T\) un arbre à niveaux finis non vides tel que tout nœud
survivant ait un descendant survivant. Alors il existe une branche
infinie. Si, en outre, chaque nœud de la sous-tour sélectionnée a exactement
un descendant compatible, la branche est unique.
\end{lemma}

\begin{proof}
L'existence résulte du lemme de König appliqué à l'arbre à branchement fini.
Pour l'unicité, deux branches distinctes auraient un premier niveau où elles
divergent, donnant deux descendants compatibles du même préfixe, contrairement à
l'hypothèse.
\end{proof}

\begin{controlbox}
Une exécution à la profondeur \(N\) ne prouve que ce qui s'est produit dans l'arbre
tronqué. L'énoncé infini exige la loi uniforme qui préserve la non-
vacuité et la compatibilité à tout niveau. Le corps du texte sépare expressément
les deux statuts et, pour la filtration nonadique, conserve son sélecteur uniforme
certifié.
\end{controlbox}

\section{Clôture et pleine puissance}

Pour \(A\subseteq X=\varprojlim X_n\), il a été prouvé
\[
 \overline A=\bigcap_n p_n^{-1}(p_n(A)).
\]
L'ensemble des suites binaires ultimement nulles est dense et propre dans
\(2^{\mathbb N}\), mais possède toutes les ombres finies. Ainsi, tout
algorithme qui ne reçoit que \((p_n(A))_n\) ne peut pas distinguer \(A\) de sa
clôture. Cet exemple bloque la fausse identité
\[
 \varprojlim\mathcal P(X_n)=\mathcal P(\varprojlim X_n).
\]

\section{Compatibilité typée des opérateurs}

La différence entre l'inclusion unitaire et le coin marqué se vérifie
directement :
\[
\begin{aligned}
 \tau_{d+1}(a\otimes I_9)
 &=\frac{\operatorname{Tr}(a)\operatorname{Tr}(I_9)}{9^{d+1}}
 =\tau_d(a),\\
 \tau_{d+1}(a\otimes p_j)
 &=\frac{\operatorname{Tr}(a)}{9^{d+1}}
 =\frac1{9}\tau_d(a).
\end{aligned}
\]
La première branche est unitaire et traciale ; la seconde conserve un coin marqué
et ne peut pas être utilisée pour inférer le facteur \(II_1\).

\section{Annulation du coût de Landauer dans l'inversion réversible de l'état}

Pour une variable complète \(X\) et une projection \(Y=q(X)\),
\[
 H(X)=H(Y)+H(X\mid Y).
\]
Si \(q\) est bijective, la fibre est un singleton et \(H(X\mid Y)=0\). Si \(q\)
est une réinitialisation du TRIT uniforme, l'information rejetée est \(\log3\) et
le coût minimal est \(k_B\Theta\log3\). Pour l'extension
\(\widehat E_N\) à \(L^1\) d'une espérance conditionnelle normale
préservant la trace \(E_N\), et sous les hypothèses de finitude déclarées,
\[
 S_\tau(\widehat E_Nh)-S_\tau(h)
 =D_\tau(h\Vert \widehat E_Nh)\ge0.
\]
Les deux formules situent le même phénomène dans des domaines différents : le
coût appartient à la contraction d'une fibre, non à l'évolution réversible
de l'état complet.

\section{Obstruction de Gödel–Turing à un décideur uniforme des multiplicités projectives}

Soit \(T_e\) la machine d'indice \(e\). Au niveau \(n\), on définit
\[
 X_n^e=\{b_n\}\cup
 \{c_n:T_e\text{ ne s'est pas arrêté avant }n\}.
\]
Chaque \(X_n^e\) est calculable. La fibre globale a deux branches si \(T_e\) ne
s'arrête pas et une si elle s'arrête. Par conséquent
\[
 U_e:\quad \#q_\infty^{-1}(v)=1
 \quad\Longleftrightarrow\quad T_e\downarrow.
\]
Un décideur uniforme pour \(U_e\) déciderait \(\HALT\). Une théorie
récursivement axiomatisable, correcte et complète pour tous les \(U_e\)
permettrait de les décider en énumérant en parallèle les preuves de \(U_e\) et de sa
négation. C'est le noyau du \emph{Théorème HMT d'incomplétude uniforme
pour les multiplicités projectives} ; sa démonstration utilise une réduction
de Gödel--Turing. L'obstruction uniforme ne nie pas les sélecteurs particuliers
construits sur des images HMT calibrées.

\begin{controlbox}[title={Portée du contrôle d'incomplétude uniforme}]
La réduction de Gödel--Turing établit l'obstruction uniforme pour les
multiplicités projectives. Elle ne nie pas les sélecteurs particuliers construits
sur des images HMT calibrées et ne transforme pas une obstruction de décision globale
en une indétermination de chaque fibre concrète.
\end{controlbox}
